% Adapted from the supplied FirstVersion.tex (Ryuya Hora template).
% The environment-specific dmjhira font declaration is omitted for portability.
\usepackage[left=2cm,right=2cm,top=2.2cm,bottom=2.2cm]{geometry}
\usepackage[utf8]{inputenc}
\usepackage[T1]{fontenc}
\usepackage{amsfonts,amsthm,amssymb,mathtools,etoolbox,stmaryrd}
\usepackage[colorlinks=true,urlcolor=blue,linkcolor=blue,citecolor=blue]{hyperref}
\usepackage{tikz,tikz-cd}
\usetikzlibrary{calc,arrows.meta,positioning,fit,backgrounds}
\usepackage{cleveref}
\usepackage{xcolor}
\usepackage{array,booktabs,longtable,tabularx}
\usepackage{enumitem}
\usepackage{microtype}
\usepackage{xurl}
\usepackage{listings}
\usepackage[style=alphabetic,sorting=nyt,maxnames=5]{biblatex}
\renewbibmacro{in:}{}
\addbibresource{CommonBiblio20240922.bib}
\addbibresource{gsh_additions.bib}

\tikzset{pullback/.style={minimum size=1.2ex,path picture={
\draw[opacity=1,black,-,#1] (-0.5ex,-0.5ex) -- (0.5ex,-0.5ex) -- (0.5ex,0.5ex);%
}}}

\theoremstyle{plain}
\newtheorem{theorem}{Theorem}[subsection]
\newtheorem{proposition}[theorem]{Proposition}
\newtheorem{lemma}[theorem]{Lemma}
\newtheorem{corollary}[theorem]{Corollary}
\newtheorem{todo}[theorem]{Todo}
\newtheorem{conjecture}[theorem]{Conjecture}
\newtheorem{fact}[theorem]{Fact}
\newtheorem{claim}{Claim}[theorem]
\crefname{claim}{Claim}{Claims}
\renewcommand{\theclaim}{\thetheorem.\alph{claim}}

\theoremstyle{definition}
\newtheorem{example}[theorem]{Example}
\newtheorem{definition}[theorem]{Definition}
\newtheorem{notation}[theorem]{Notation}
\newtheorem{puzzle}[theorem]{Puzzle}
\newtheorem{idea}[theorem]{Idea}
\newtheorem{question}[theorem]{Question}
\newtheorem{exercise}[theorem]{Exercise}
\newtheorem{protocol}[theorem]{Protocol}

\theoremstyle{remark}
\newtheorem{remark}[theorem]{Remark}
\newtheorem{warning}[theorem]{Warning}

\newcommand{\dq}[1]{``#1''}
\newcommand{\memo}[1]{\textcolor{red}{memo: #1}}
\newcommand{\invmemo}[1]{}
\newcommand{\horamemo}[1]{\textcolor{green!60!black}{hora: #1}}
\newcommand{\para}[1]{\paragraph{\textbf{#1}}}
\newcommand{\N}{\mathbb{N}}
\newcommand{\Z}{\mathbb{Z}}
\newcommand{\Q}{\mathbb{Q}}
\newcommand{\R}{\mathbb{R}}
\newcommand{\C}{\mathcal{C}}
\newcommand{\D}{\mathcal{D}}
\newcommand{\E}{\mathcal{E}}
\newcommand{\F}{\mathbf{F}}
\renewcommand{\L}{\mathcal{L}}
\newcommand{\id}{\mathrm{id}}
\newcommand{\op}{\mathrm{op}}
\newcommand{\ob}{\mathrm{ob}}
\newcommand{\Set}{\mathbf{Set}}
\newcommand{\FinSet}{\mathbf{FinSet}}
\newcommand{\Func}[2]{[#1,#2]}
\newcommand{\abs}[1]{\left|#1\right|}
\newcommand{\demph}[1]{\textit{#1}}
\newcommand{\Pow}{\mathcal{P}}
\newcommand{\Words}[1]{#1^{\ast}}
\newcommand{\eps}{\varepsilon}
\newcommand{\sh}{\operatorname{sh}}
\newcommand{\Syn}{\operatorname{Syn}}
\newcommand{\Aper}{\mathbf{A}}
\newcommand{\Afour}{A_{4}}
\newcommand{\Afive}{A_{5}}
\newcommand{\Dicthree}{\operatorname{Dic}_{3}}
\newcommand{\heightone}{\mathsf{H}_{1}}
\newcommand{\code}[1]{\texttt{\detokenize{#1}}}
\newcommand{\status}[1]{\textsf{#1}}

\definecolor{codegray}{gray}{0.97}
\lstdefinestyle{workshop}{
  basicstyle=\ttfamily\small,
  backgroundcolor=\color{codegray},
  frame=single,
  framerule=0.2pt,
  columns=fullflexible,
  keepspaces=true,
  breaklines=true,
  showstringspaces=false,
  upquote=true,
  literate={α}{{$\alpha$}}1
}
\lstset{style=workshop}

\setlist[itemize]{topsep=3pt,itemsep=2pt,parsep=0pt}
\setlist[enumerate]{topsep=3pt,itemsep=2pt,parsep=0pt}
\setcounter{tocdepth}{2}
\setcounter{secnumdepth}{3}
